\chapter{Functions: inputs, outputs, and transformations}\label{ch:functions}
\begin{readergoal}Distinguish relations from functions; work with domain, codomain, image, inverse, and composition.\end{readergoal}
\section{A function is a special relation}
A function $f:A\to B$ assigns to each $a\in A$ exactly one value $f(a)\in B$. It may assign the same output to many inputs. The set $A$ is its \emph{domain}, $B$ its \emph{codomain}, and $f(A)=\{f(a):a\in A\}$ its image or range. A relation that assigns two distinct outputs to the same input is not a function. These distinctions matter when saying ``the neural signal determines the behavior'': an observed relation between activity and behavior is not automatically a deterministic function.

A function is \emph{injective} when different inputs have different outputs, \emph{surjective} when every point in the codomain occurs as an output, and \emph{bijective} when both conditions hold. Only a bijection $f:A\to B$ has a genuine inverse function $f^{-1}:B\to A$. Every relation has an inverse relation, but an arbitrary function does not have an inverse function.

\section{Compositions and transformations}
Given $f:A\to B$ and $g:B\to C$, their composition is $g\circ f:A\to C$, where $(g\circ f)(a)=g(f(a))$. Composition generally does not commute. In electrophysiology, one might take raw voltage $x$, apply a filter $F$, and then a nonlinear amplitude transform $H$: $H(F(x))$. Reversing the order may be meaningless or produce a different result.

\section{Sequences are functions}
A sequence $(x_n)_{n\in\mathbb N}$ is a function from an index set of natural numbers into a set of values. A finite sampled signal is a function $x:\{0,\ldots,N-1\}\to\mathbb R$. A continuously indexed signal is a function $x:I\to\mathbb R$ for some time interval $I$. A multichannel signal is $x:I\to\mathbb R^C$. The difference between continuous and sampled time is not a notational accident: sampling introduces an observation process and potentially destroys information.

\section{Functions of functions}
A \emph{functional} maps functions to numbers. For example, energy over an interval can be defined as $E[x]=\int_I |x(t)|^2\,dt$. A \emph{linear operator} maps functions to functions while respecting addition and scalar multiplication. Filtering, differentiation, and Fourier transformation are operators under suitable assumptions. These ideas will become precise after we construct number systems, vector spaces, and limits.

\begin{workedexample}[A map that loses information]
The square map $q:\mathbb R\to\mathbb R$, $q(x)=x^2$, is not injective: $q(-2)=q(2)=4$. Its inverse relation associates $4$ with both $2$ and $-2$. Restricting the domain to $[0,\infty)$ and codomain to $[0,\infty)$ makes $q$ bijective and its inverse is the nonnegative square root. Signal power similarly discards phase and sign information: it cannot reconstruct the original signal by itself.
\end{workedexample}
% --- expanded edition (October 2026) ---
\section{Images and preimages of sets}
For $f:A\to B$ and $S\subseteq A$, the \emph{image} of $S$ is $f(S)=\{f(a):a\in S\}$; for $T\subseteq B$, the \emph{preimage} is $f^{-1}(T)=\{a\in A:f(a)\in T\}$. The preimage is defined for every function, whether or not $f$ is invertible; the notation $f^{-1}(T)$ does not presuppose an inverse function.

\begin{proposition}[Preimages respect set operations]
For $f:A\to B$ and subsets $T,T_i\subseteq B$:
\[
f^{-1}\Bigl(\bigcup_iT_i\Bigr)=\bigcup_if^{-1}(T_i),\quad
f^{-1}\Bigl(\bigcap_iT_i\Bigr)=\bigcap_if^{-1}(T_i),\quad
f^{-1}(B\setminus T)=A\setminus f^{-1}(T).
\]
Images respect unions, $f(\bigcup_iS_i)=\bigcup_if(S_i)$, but in general only $f(S_1\cap S_2)\subseteq f(S_1)\cap f(S_2)$.
\end{proposition}
\begin{proof}
$a\in f^{-1}(\bigcap_iT_i)$ iff $f(a)\in T_i$ for every $i$ iff $a\in f^{-1}(T_i)$ for every $i$. The union and complement cases are identical in form. For images and intersections, if $b=f(a)$ with $a\in S_1\cap S_2$ then $b\in f(S_1)$ and $b\in f(S_2)$. Equality can fail: with $f(x)=x^2$, $S_1=\{-1\}$, $S_2=\{1\}$, we have $f(S_1\cap S_2)=\varnothing$ but $f(S_1)\cap f(S_2)=\{1\}$.
\end{proof}
This asymmetry explains why measurability (Part~V) is defined through preimages: a function is measurable when preimages of measurable sets are measurable, precisely because preimages commute with all set operations.

\section{Left and right inverses}
A function $g:B\to A$ is a \emph{left inverse} of $f:A\to B$ if $g\circ f=\mathrm{id}_A$, and a \emph{right inverse} if $f\circ g=\mathrm{id}_B$.

\begin{theorem}
Let $A$ be nonempty and $f:A\to B$.
\begin{enumerate}[leftmargin=2em]
\item $f$ has a left inverse if and only if $f$ is injective.
\item $f$ has a right inverse if and only if $f$ is surjective (using the axiom of choice for infinite $B$).
\item If $f$ has both a left inverse $g$ and a right inverse $h$, then $g=h$, and this common function is the unique inverse $f^{-1}$.
\end{enumerate}
\end{theorem}
\begin{proof}
(1) If $g\circ f=\mathrm{id}$ and $f(a)=f(a')$, then $a=g(f(a))=g(f(a'))=a'$. Conversely, if $f$ is injective, fix $a_0\in A$ and define $g(b)$ to be the unique $a$ with $f(a)=b$ when $b\in f(A)$, and $a_0$ otherwise. (2) If $f\circ h=\mathrm{id}$, every $b=f(h(b))$ lies in the image. Conversely, if $f$ is surjective, each preimage $f^{-1}(\{b\})$ is nonempty; choosing one element from each defines $h$. (3) $g=g\circ\mathrm{id}_B=g\circ(f\circ h)=(g\circ f)\circ h=\mathrm{id}_A\circ h=h$.
\end{proof}

A left inverse is a \emph{decoder} that recovers the input whenever the output came from the encoder. A right inverse is a \emph{generator} that produces, for each target output, some input achieving it. In neural decoding, a stimulus-to-response map that is not injective (two stimuli evoke the same response) has no perfect decoder; in neural control, a response space larger than the set of achievable responses has no right inverse.

\section{Composition, injectivity, and surjectivity}
\begin{proposition}
Let $f:A\to B$ and $g:B\to C$.
(i) If $f$ and $g$ are injective, so is $g\circ f$.
(ii) If $f$ and $g$ are surjective, so is $g\circ f$.
(iii) If $g\circ f$ is injective, then $f$ is injective.
(iv) If $g\circ f$ is surjective, then $g$ is surjective.
\end{proposition}
\begin{proof}
(i) $g(f(a))=g(f(a'))$ gives $f(a)=f(a')$ and then $a=a'$. (ii) For $c\in C$ choose $b$ with $g(b)=c$, then $a$ with $f(a)=b$. (iii) If $f(a)=f(a')$ then $g(f(a))=g(f(a'))$, so $a=a'$. (iv) Every $c=g(f(a))$ is a value of $g$.
\end{proof}
Statement (iii) has an important consequence for processing pipelines: if any early stage $f$ discards information (is not injective), no later stage $g$ can restore it. Rectification, thresholding, binning, and taking power are all non-injective stages.

\section{Cardinality and countability}
Sets $A$ and $B$ have the \emph{same cardinality}, $|A|=|B|$, if there is a bijection between them. A set is \emph{countable} if it is finite or has the cardinality of $\mathbb N_0$. The integers are countable via $0,1,-1,2,-2,\ldots$, and so are the rationals and every finite product of countable sets.

\begin{theorem}[Cantor]
For every set $A$ there is no surjection $A\to\mathcal P(A)$. In particular $|A|<|\mathcal P(A)|$.
\end{theorem}
\begin{proof}
Let $f:A\to\mathcal P(A)$ and set $D=\{a\in A:a\notin f(a)\}$. If $D=f(d)$ for some $d$, then $d\in D\iff d\notin f(d)=D$, a contradiction. Hence $D$ is not in the image of $f$.
\end{proof}
Applied to $A=\mathbb N_0$, Cantor's theorem shows that the set of infinite binary sequences (equivalently, the set of all possible infinite spike rasters at fixed resolution) is uncountable. This is one reason why statements about ``all possible spike trains'' require measure theory rather than counting.
\begin{exercises}
\item Decide whether $f(x)=x^3$ from $\mathbb R$ to $\mathbb R$ is bijective.
\item State domain and codomain for a histogram of spike counts in 10-ms bins.
\item Find $f\circ g$ and $g\circ f$ for $f(x)=x+1$ and $g(x)=x^2$.
\item Explain why a spectrogram generally does not contain sufficient information to uniquely recover a waveform without additional phase assumptions.
\item Prove that $f(f^{-1}(T))\subseteq T$ with equality when $f$ is surjective, and $S\subseteq f^{-1}(f(S))$ with equality when $f$ is injective.
\item Give a function with a left inverse but no right inverse, and describe all its left inverses.
\item Show that $g\circ f$ injective does not imply $g$ injective.
\item Construct an explicit bijection $\mathbb N_0\to\mathbb Z$.
\item A spike detector maps a 30\,kHz voltage trace to a list of threshold-crossing times. Identify the domain, codomain, and why the map is not injective.
\end{exercises}
